% Chapter 1

\chapter{Introducción general}

\label{Chapter1}
\label{IntroGeneral}

En este capítulo se presenta el contexto del sistema desarrollado para heladeras comerciales, la problemática que motivó el trabajo y los principales antecedentes relacionados. Finalmente, se establecen los objetivos y el alcance del desarrollo.

%----------------------------------------------------------------------------------------

\section{Sistema Sensify para heladeras comerciales}

Sensify S. A. S. \cite{sensify:web} es una empresa dedicada al análisis de datos que trabaja junto con empresas de alimentos y bebidas de América Latina para convertir heladeras comerciales en fuentes de información que permitan gestionar sus ventas de manera inteligente y sostenible. Uno de sus productos centrales es el controlador N150 \cite{sensify:n150}, un dispositivo diseñado para integrarse en heladeras comerciales con el propósito de supervisar su funcionamiento y transmitir información relevante hacia plataformas remotas mediante conectividad celular (2G y 4G).

El controlador integra en una única plataforma funciones de control, monitoreo y conectividad \cite{sensify:n150}. Esta integración permite disponer de información sobre el estado operativo de las heladeras instaladas y facilita las tareas de diagnóstico y mantenimiento. Además, la información obtenida contribuye a la gestión técnica y operativa del parque de equipos conectados.

De acuerdo con datos internos de Sensify correspondientes a septiembre de 2026, la plataforma contaba con un parque aproximado de 70\,000 heladeras conectadas, frente a las 8\,000 unidades registradas en 2023. Para esa fecha, 8 fabricantes de heladeras habían integrado la tecnología de Sensify en sus equipos y la solución se encontraba presente en 10 países.

{\textcolor{red}{Aclaracion para la correcion: Estoy en proceso de que la empresa me proporcione imagenes para ilustrar esta seccion.}}

\subsection{El controlador N150}

Desde el punto de vista arquitectónico, el controlador se basa en la interacción de distintos componentes especializados. Incorpora el PIC24FJ64GU205T \cite{microchip:pic24fj64gu205}, un microcontrolador encargado de ejecutar la lógica de control eléctrico y mecánico del sistema de refrigeración. También integra un microcontrolador ESP32 \cite{espressif:esp32}, responsable de la conectividad y del procesamiento asociado a las comunicaciones. Para el acceso a redes celulares mediante comandos AT, cuenta con un módulo Quectel BG95 \cite{quectel:bg95}. Además, dispone de una memoria externa para almacenar información y respaldar datos cuando las condiciones de conectividad no permiten su transmisión inmediata.

Esta estructura proporciona capacidades de telemetría, parametrización remota, actualización de firmware y geolocalización de las heladeras instaladas. Con estas funciones, el N150 opera no solo como controlador del equipo de frío, sino también como nodo inteligente dentro de una red de activos conectados.

%----------------------------------------------------------------------------------------

\section{Contexto y motivación}

La motivación del trabajo surgió de una limitación observada durante las tareas de configuración y soporte técnico del N150. Si bien la conectividad celular permitía atender las necesidades de telemetría y transmisión remota de datos, presentaba limitaciones para las tareas que requerían una interacción inmediata con el equipo. Entre estas tareas se encontraban la activación inicial, la modificación de parámetros, el diagnóstico técnico y la actualización de firmware.

Ante esta situación, se identificó la necesidad de disponer de un mecanismo de comunicación local que permitiera efectuar dichas operaciones de forma directa, rápida y confiable, sin depender de la cobertura ni de la estabilidad de la red celular. Esta necesidad adquiría especial relevancia durante la instalación, la validación en planta y el soporte técnico en campo, debido a que la disponibilidad de un canal de acceso inmediato condicionaba la eficiencia de cada intervención.

La incorporación de un canal local también ofrecía una oportunidad de mejora operativa. Por este motivo, se desarrolló un sistema de comunicación basado en Bluetooth para interactuar con el controlador mediante un dispositivo móvil de uso común. El sistema tuvo como finalidad reducir la necesidad de herramientas específicas, simplificar las tareas de mantenimiento y ampliar las capacidades de soporte sobre los equipos instalados.

Desde el punto de vista formativo, el desarrollo permitió abordar un problema real de ingeniería mediante la integración de conceptos de sistemas embebidos, comunicaciones inalámbricas, arquitectura de software y diseño de interfaces de usuario.

%----------------------------------------------------------------------------------------

\section{Estado del arte}

Las soluciones conectadas aplicadas a la refrigeración comercial adoptan distintos enfoques según la conectividad empleada y las funciones disponibles. Algunas se concentran en la adquisición de variables y el monitoreo remoto, mientras que otras incorporan funciones de configuración, diagnóstico o control. Para analizar estos enfoques se consideraron soluciones comerciales y antecedentes académicos relacionados con la supervisión de equipos refrigerados y el acceso local a sistemas embebidos.

Entre las soluciones comerciales orientadas al monitoreo se encuentran Swift Sensors \cite{SwiftSensors:web} y Monnit \cite{Monnit:web}. Ambas ofrecen sensores inalámbricos, plataformas de supervisión remota y generación de alertas. Swift Sensors utiliza Bluetooth entre los sensores y una puerta de enlace que transmite la información hacia una plataforma remota. Monnit dispone de sensores, puertas de enlace y herramientas de gestión para distintas variables ambientales y operativas. La documentación consultada de estas plataformas se concentra en la adquisición, el registro y la visualización de datos, pero no describe funciones de configuración profunda ni de actualización del software embebido del controlador refrigerado.

Otro enfoque consiste en proporcionar acceso local mediante tecnologías de corto alcance. El controlador X35P de Coel admite, de manera opcional, la programación de parámetros mediante comunicación de campo cercano (NFC, del inglés, \textit{Near Field Communication}) \cite{Coel:web}. Este mecanismo simplifica la configuración en el lugar de instalación y no requiere una red de datos externa. Sin embargo, exige proximidad física y no está orientado al monitoreo remoto continuo.

Las soluciones basadas en conectividad celular evitan depender de la red local del establecimiento. AoFrio ofrece controladores conectados y la plataforma AoFrio iQ, que proporciona monitoreo, diagnóstico y control remoto de equipos de refrigeración comercial \cite{AoFrio:web,AoFrioIQ:web}. Este esquema resulta apropiado para administrar flotas distribuidas, aunque su funcionamiento remoto depende de la cobertura y de la calidad de la señal celular. En ubicaciones con servicio deficiente, la comunicación puede demorarse o interrumpirse.

La tabla \ref{tab:comparacion_estado_arte} resume las características relevantes de las soluciones analizadas. La comparación distingue entre monitoreo, configuración o control y disponibilidad de acceso local, debido a que estas funciones responden a necesidades operativas diferentes.

\begin{table}[h]
	\centering
	\caption[Comparación de soluciones]{Comparación de las soluciones comerciales y académicas analizadas.}
	\small
	\begin{tabular}{p{2.3cm} p{2.6cm} p{3.8cm} p{3.2cm}}
		\toprule
		\textbf{Solución} & \textbf{Conectividad} & \textbf{Funciones principales} & \textbf{Limitación relevante} \\
		\midrule
		Swift Sensors & Bluetooth y puerta de enlace & Monitoreo remoto, registro y alertas & No documenta configuración profunda del controlador \\
		Monnit & Red inalámbrica y puerta de enlace & Monitoreo remoto, registro y alertas & No documenta actualización de firmware del controlador \\
		Coel X35P & NFC y RS485 & Configuración local de parámetros y control & Requiere proximidad y no brinda monitoreo remoto continuo \\
		AoFrio iQ & Celular & Monitoreo, diagnóstico y control remoto & Depende de la cobertura celular \\
		\bottomrule
		\hline
	\end{tabular}
	\label{tab:comparacion_estado_arte}
\end{table}

En el N150, la memoria externa permite conservar la información generada durante los períodos sin conectividad y transmitirla cuando se restablece el enlace. Este mecanismo resulta adecuado para la transmisión diferida de telemetría, pero no proporciona el acceso inmediato requerido durante la instalación y el soporte técnico. Entre los antecedentes analizados no se identificó una solución que reuniera la conectividad celular existente en el N150 con un canal local destinado al diagnóstico, la configuración y la actualización de firmware. Esta diferencia fundamentó la incorporación de Bluetooth como canal complementario, sin reemplazar el enlace celular utilizado para la operación remota.

%----------------------------------------------------------------------------------------

\section{Objetivos del trabajo}

El objetivo general del trabajo fue diseñar e implementar un sistema de comunicación local basado en Bluetooth para el controlador Sensify N150. El sistema debía complementar la conectividad remota existente y permitir la ejecución de tareas de soporte técnico, configuración y mantenimiento sin depender de la disponibilidad de la red celular.

Para alcanzar el objetivo general se definieron los siguientes objetivos específicos:

\begin{itemize}
	\item Incorporar en el firmware del ESP32 un módulo que permitiera establecer una comunicación inalámbrica local con dispositivos móviles.
	\item Permitir la consulta de información de estado y diagnóstico del controlador mediante el nuevo canal de comunicación.
	\item Permitir la modificación de parámetros operativos y la activación inicial del equipo.
	\item Implementar un mecanismo para actualizar localmente el firmware mediante Bluetooth.
	\item Integrar la nueva funcionalidad en la arquitectura existente del N150 sin interferir con los mecanismos de conectividad previamente implementados.
	\item Administrar la memoria dinámica y los recursos compartidos utilizados por los distintos módulos de software.
	\item Verificar mediante pruebas funcionales la comunicación, las operaciones de soporte implementadas y la estabilidad general del sistema.
\end{itemize}

Estos objetivos establecieron resultados verificables para la comunicación local, las funciones disponibles y la integración con el sistema existente. Los criterios y procedimientos utilizados para su evaluación se presentan en los capítulos posteriores.
